\section{Discussion}
\label{sec:disc}

The remaining-exposure construction establishes a population-dependent floor
against arbitrary admitted feedback. The product in that floor is a function of
the whole configuration, not an independence approximation during treatment, and
its founder dependence survives the coupling that a common controller
introduces. The amplitude construction is a different certificate with a
different conclusion: below a critical actuator strength the survival floor does
not decay at all, so no budget and no schedule can help. Separating the two is
what makes the memory-size question well posed, and what turns a diagnostic
about one constant action into a statement about every policy.

Three concrete conclusions follow for the molecular source. First, the
obstruction that scales badly with the size of the inherited state is
\emph{amplitude}, not exposure: between two and eight sites the certified
critical erasure rises by a factor $4.86$ while the necessary exposure rises by
$22\%$. Second, the amplitude used in the two-site worked example provably fails
at five to eight sites, against every predictable feedback policy, at infinite
budget and infinite time. Third, the failure is a property of that actuator:
added death on protected states has a threshold bounded by $b-d_{\mathrm{prot}}$
for every $N$ by an elementary total-count argument, delivers a dimension-free
sufficient exposure $1.45\log(n/\delta)$, and attains the exposure floor to
within a factor $1.533$ at $n=1$, $\delta=.01$. That last point also supplies
the sharpness that Proposition~\ref{prop:sharp} could only give in a degenerate
comparison source.

What remains open is quantitative and specific. For the eraser at two sites the
certified bracket at $T=100$ is $4.7257\le B^*\le11.6$, and
Remark~\ref{rem:optcert} shows the lower endpoint is within $0.54\%$ of what
certificate optimization alone can deliver, so the residual factor of about
$2.45$ is not a defect of the chosen supersolution. Closing it requires either a
richer population barrier or a better policy, and the natural next experiment is
to certify a small phase family selected by backward sensitivity of the
generating-function flow and compare it with a population-budget barrier that
does not factor through a single vector $q$. Whether $e_c(N)$ remains bounded as
$N\to\infty$ is open; the exploratory evidence of Remark~\ref{rem:bigN} is
consistent with either answer, and a proof in either direction would say
something structural about how fast inherited memory can be defeated by a
single reaction-rate actuator.

Uniform constants in site count for the eraser, density-dependent populations,
empirical target engagement, a global solution of the associated dynamic
programming equation, and full stochastic formalization all remain outside this
paper. They are explicit future scopes and not hidden assumptions: the
formalization boundary is stated in Appendix~\ref{app:boundary}, the evidence
type of every displayed quantity is labelled, and the exact vectors behind every
certificate are released with the source package.
